\section{An integral class that persists at every scale}
\label{sec:persistence}

We construct a class $\beta\in\KX_1(B)$ whose images at the two ends of
$X$ have an infinite-order difference. The obstruction to identifying
these images will be integral: at any fixed Rips scale, an appropriate
$K$-cohomology class becomes divisible by larger and larger powers of two
over the middle graphs $Y_j$.

All $K$-groups in this section are complex and integral. For a finite
complex $M$, we write $K^*(M)$ for topological $K$-theory and
$K_*(M)=\KK_*(C(M),\C)$ for analytic $K$-homology. The pairing
\[
 \langle\ ,\ \rangle:K^1(M)\times K_1(M)\longrightarrow\Z
\]
is the Kasparov index pairing. In particular, for a continuous map
$h:M\to N$ between finite complexes,
\begin{equation}
 \langle\eta,h_*\xi\rangle=\langle h^*\eta,\xi\rangle.
 \label{top:pairing-naturality}
\end{equation}
We use the usual Mayer--Vietoris sequence for a finite complex covered
by two subcomplexes; see \cite{HigsonRoe00} for these $K$-homology
constructions.

\subsection{The classes on the two ends}

Let
\[
 T_j=(\R/l_j\Z)^3.
\]
Give $T_j$ its unit cubical subdivision and triangulate the cubes
periodically. The vertices are $B_j$, and the vertices of every cube
have graph distance at most $3$. The resulting triangulated torus is
therefore a subcomplex of $P_3(B_j)$; write
$\kappa_j:T_j\hookrightarrow P_3(B_j)$ for its inclusion.

Let $u_1,u_2,u_3\in H^1(T_j;\Z)$ be the positively oriented coordinate
classes. Choose a line bundle $\mathcal L_j$ pulled back from the first
two circles with
\[
 c_1(\mathcal L_j)=u_1\smile u_2,
\]
and let $e_3\in K^1(T_j)$ be the pullback of the positive generator of
$K^1(S^1)$ under projection to the third circle. Define
\begin{equation}
 \theta_j=([\mathcal L_j]-1)e_3\in K^1(T_j).
 \label{top:theta}
\end{equation}
This class is primitive. Indeed, the exterior-algebra description of
$K^*((S^1)^3)$ identifies $[\mathcal L_j]-1$ with a generator of the
reduced $K^0$ of the first two circles, and $\theta_j$ with the product
of the three circle generators. Here the first assertion also follows
from the first Chern class, which identifies
$\widetilde K^0((S^1)^2)$ with $H^2((S^1)^2;\Z)$.
The K\"unneth and universal coefficient theorems therefore give a class
\begin{equation}
 \tau_j\in K_1(T_j),\qquad
 \langle\theta_j,\tau_j\rangle=1.
 \label{top:tau}
\end{equation}
More explicitly, $K^*(T_j)$ is free abelian, so the universal
coefficient theorem identifies $K_1(T_j)$ with
$\operatorname{Hom}(K^1(T_j),\Z)$, and a primitive element admits an
integer-valued functional taking value $1$ on it; see
\cite[Theorems 23.1.1 and 23.1.3]{Blackadar98}.

The component separation gives
$P_3(B)=\bigsqcup_jP_3(B_j)$. Countable additivity in the first variable
of $\KK$ gives
\begin{equation}
 \KK_1(C_0(P_3(B)),\C)
 \cong\prod_{j\geq1}K_1(P_3(B_j)).
 \label{top:kk-product}
\end{equation}
This is a product, so the entire sequence
$((\kappa_j)_*\tau_j)_j$ determines a class $\beta_3$ in the group on
the left; see \cite[Theorem 19.7.1]{Blackadar98}. Write $\beta$ for its
image in $\KX_1(B)$ and set
\begin{equation}
 \alpha_3=(i_{+,3})_*\beta_3-(i_{-,3})_*\beta_3,
 \qquad
 \alpha=[\alpha_3]\in\KX_1(X).
 \label{top:alpha}
\end{equation}
Thus $\alpha=(i_+)_*\beta-(i_-)_*\beta$. We will prove that every
nonzero integer multiple of $\alpha_3$ remains nonzero at every later
Rips scale.

\subsection{Retractions to the tori}

Fix an integer $r\geq3$. From now on, thresholds on $j$ may depend on
$r$, but not on a simplex or a vertex of the corresponding component.

\begin{lemma}
\label{top:interpolation}
If $l_j>4r$, there is a continuous map
\[
 p_{B,j}:P_r(B_j)\longrightarrow T_j
\]
whose restriction to $T_j\subseteq P_3(B_j)\subseteq P_r(B_j)$ is the
identity. The graph projections define continuous maps
\[
 p_{Y,j}:P_r(Y_j)\longrightarrow T_j,
 \qquad
 p_{X,j}:P_r(X_j)\longrightarrow T_j
\]
by composition with $p_{B,j}$.
\end{lemma}

\begin{proof}
For vertices $u,v$ of a Rips simplex, each coordinate difference has a
unique representative $\delta_i(u,v)\in\Z$ with
$|\delta_i(u,v)|\leq r$. On a triangle $(u,v,w)$, the integer
\[
 \delta_i(u,v)+\delta_i(v,w)-\delta_i(u,w)
\]
is a multiple of $l_j$ and has absolute value at most $3r<l_j$; it
therefore vanishes. Consequently the vertices of each simplex admit
lifts to $\Z^3$ realizing all these differences. The lifts are unique
up to a common element of $l_j\Z^3$. Affine interpolation, followed by
reduction modulo $l_j\Z^3$, gives a map on the simplex independent of
that choice. The maps agree on faces and define $p_{B,j}$.

On each simplex of the triangulated torus, this interpolation is its
original parametrization, so the restriction is the identity. The
projections $Y_j\to B_j$ and $X_j\to B_j$ are $1$-Lipschitz by
Lemma~\ref{geom:basic}; hence they induce simplicial maps of the
corresponding Rips complexes.
\end{proof}

The maps in Lemma~\ref{top:interpolation} allow us to compare the same
torus class at the ends and in the middle. At either end its pairing
with $\tau_j$ is $1$. In the middle, a congruence calculation forces
divisibility.

\subsection{The congruence calculation}

\begin{lemma}
\label{top:chern-divisibility}
For every fixed integer $r\geq3$ and all sufficiently large $j$,
\begin{equation}
 p_{Y,j}^*(u_1\smile u_2)
 \in 2^j H^2(P_r(Y_j);\Z).
 \label{top:chern-divisibility-eq}
\end{equation}
Equivalently, there is a line bundle $\mathcal L'_{j,r}$ on
$P_r(Y_j)$ such that
\begin{equation}
 p_{Y,j}^*\mathcal L_j
 \cong(\mathcal L'_{j,r})^{\otimes 2^j}.
 \label{top:line-root}
\end{equation}
\end{lemma}

\begin{proof}
Take $j$ large enough that $l_j>4r$ and reduction
$\Gamma\to Q_{2j}$ is injective on the word ball of radius $2r$;
the latter condition is available from Lemma~\ref{geom:quotients}.
In this proof write $m=2^j$, suppress $j$ in $U_j,V_j$, and abbreviate
$g_{j,k}$ to $g_k$.

We first describe integral cocycles representing
$p_{Y,j}^*u_1,p_{Y,j}^*u_2$. For an oriented Rips edge $(x,y)$, let
$v_x,v_y\in B_j$ be its projected endpoints and choose their coordinate
representatives in $\{0,\ldots,l_j-1\}$. Define its first two seam
numbers $a_{xy},b_{xy}\in\Z$ by
\[
 \delta_1(v_x,v_y)=(v_y)_1-(v_x)_1+l_j a_{xy},
 \qquad
 \delta_2(v_x,v_y)=(v_y)_2-(v_x)_2+l_j b_{xy}.
\]
The triangle identities in the proof of Lemma~\ref{top:interpolation}
show that $a,b$ are cocycles. The seam numbers record the winding of
the interpolated coordinate maps, so
$[a]=p_{Y,j}^*u_1$ and $[b]=p_{Y,j}^*u_2$.

To prove divisibility by $m$, it suffices to show that $a\smile b$
becomes a coboundary modulo $m$. We will construct a
$\Z/m\Z$-valued $1$-cochain $c$ with $\delta c=-a\smile b$ by
comparing edge transports after lifting from $Q_{2j}$ to $Q_{3j}$.

Order the vertices of $P_r(Y_j)$. For every edge $x<y$, choose a graph
path from $x$ to $y$ of length at most $r$, and write
$x=(v_x,q_x)$ and $y=(v_y,q_y)$. The left Cayley steps along this path
give a word $w_{xy}\in\Gamma$ of length at most $r$. Left
multiplications commute with all horizontal right multiplications.
Moreover, $U,V$ commute modulo $m^2$, and the total seam numbers of the
projected path are $a_{xy},b_{xy}$: its coordinate displacement is
bounded by $r$ and therefore is the unique small lift used above.
It follows that
\begin{equation}
 q_y=\overline w_{xy}\,q_x\,g_{2j}(a_{xy},b_{xy})
 \quad\text{in }Q_{2j}.
 \label{top:edge-transport}
\end{equation}
Here a bar denotes reduction modulo $m^2$.

On an ordered triangle with vertices $x_0<x_1<x_2$, cocycle additivity
and commutation modulo $m^2$ give
\[
 g_{2j}(a_{01},b_{01})g_{2j}(a_{12},b_{12})
 =g_{2j}(a_{02},b_{02}).
\]
Substitution in \eqref{top:edge-transport} and cancellation give
$\overline w_{02}=\overline w_{12}\overline w_{01}$. Both
$w_{02}$ and $w_{12}w_{01}$ lie in the word ball of radius $2r$, so
our injectivity assumption yields the equality in $\Gamma$:
\begin{equation}
 w_{02}=w_{12}w_{01}.
 \label{top:word-triangle}
\end{equation}

We now lift this calculation from modulus $m^2$ to modulus $m^3$.
Every element of
\[
 N=\ker(Q_{3j}\longrightarrow Q_{2j})
\]
has the form $1+m^2D$ modulo $m^3$, with $D$ determined modulo $m$.
Multiplication adds the matrices $D$, since $m^4$ is divisible by
$m^3$. Thus $N$ is abelian and
\[
 d:N\longrightarrow\Z/m\Z,
 \qquad d(1+m^2D)=D_{13}\pmod m
\]
is a homomorphism. Every element of $N$ commutes with every matrix
that is $1$ modulo $m$, in particular with
\[
 \widetilde g(a,b)=V^bU^a\pmod{m^3}.
\]
Direct multiplication of the elementary matrices gives
\begin{equation}
 \begin{split}
 C&=UVU^{-1}V^{-1}=1+m^2E_{13}\pmod{m^3},
       \qquad d(C)=1,\\
 \widetilde g(a,b)\widetilde g(a',b')
   &=\widetilde g(a+a',b+b')C^{ab'}.
 \end{split}
 \label{top:central-extension}
\end{equation}

Choose a lift $\widetilde q_x\in Q_{3j}$ of each $q_x$, and denote the
reduction of $w\in\Gamma$ to $Q_{3j}$ by $\widetilde w$.
Equation~\eqref{top:edge-transport} defines $n_{xy}\in N$ by
\begin{equation}
 \widetilde q_y=
 \widetilde w_{xy}\widetilde q_x
 \widetilde g(a_{xy},b_{xy})n_{xy}.
 \label{top:lifted-edge}
\end{equation}
For an ordered triangle, abbreviate
$\widetilde g(a_{ab},b_{ab})$ to $\widetilde g_{ab}$. Substituting
twice in \eqref{top:lifted-edge} gives
\begin{align*}
 \widetilde q_2
 &=\widetilde w_{12}\widetilde w_{01}\widetilde q_0
       \widetilde g_{01}n_{01}\widetilde g_{12}n_{12}\\
 &=\widetilde w_{02}\widetilde q_0
       \widetilde g_{01}\widetilde g_{12}n_{01}n_{12}\\
 &=\widetilde w_{02}\widetilde q_0\widetilde g_{02}
       C^{a_{01}b_{12}}n_{01}n_{12}.
\end{align*}
The second equality moves $n_{01}$ only past $\widetilde g_{12}$,
which it commutes with; it uses no commutation with
$\widetilde q_0$ or with a Cayley word. Comparing with the direct edge
$02$ therefore gives
\[
 n_{02}=C^{a_{01}b_{12}}n_{01}n_{12}.
\]
Apply $d$ and put $c_{xy}=d(n_{xy})$. With the usual simplicial
coboundary convention, this identity becomes
\[
 (a\smile b)_{012}=a_{01}b_{12}
   =c_{02}-c_{01}-c_{12}=-(\delta c)_{012}
   \quad\text{in }\Z/m\Z.
\]
Hence the reduction of $[a\smile b]$ modulo $m$ is zero. The
cohomology sequence of
\[
 0\longrightarrow\Z\xrightarrow{\,m\,}\Z
   \longrightarrow\Z/m\Z\longrightarrow0
\]
identifies the kernel of reduction on $H^2$ with the image of
multiplication by $m$. This proves
\eqref{top:chern-divisibility-eq}, without a torsion-freeness
assumption on $H^2(P_r(Y_j);\Z)$. Finally, line bundles on a finite
complex are classified by their integral first Chern class, and the
first Chern class of a tensor power is the corresponding integer
multiple. This proves \eqref{top:line-root}.
\end{proof}

\subsection{From line-bundle roots to divisibility in \texorpdfstring{$K^1$}{K1}}

For fixed $r$, bounded geometry supplies an integer $D_r\geq1$ such
that
\[
 \dim P_r(Y_j)\leq D_r\qquad(j\geq1).
\]
Indeed, the vertices of a Rips simplex all belong to the radius-$r$
ball about any one of them, and these balls have uniformly bounded
cardinality. The dimension bound converts the line-bundle roots into
integral divisibility in $K$-theory.

\begin{lemma}
\label{top:k-divisibility}
For every fixed integer $r\geq3$ and all sufficiently large $j$,
\begin{equation}
 p_{Y,j}^*\theta_j\in
 2^{\,j-\lfloor\log_2D_r\rfloor}K^1(P_r(Y_j)).
 \label{top:k-divisibility-eq}
\end{equation}
\end{lemma}

\begin{proof}
Take $j$ satisfying Lemma~\ref{top:chern-divisibility} and $2^j>D_r$,
and put
$z=[\mathcal L'_{j,r}]-1\in K^0(P_r(Y_j))$.
The classifying map of this line bundle can be taken to land in
$\mathbb{CP}^{D_r}$ by cellular approximation. The relation
$([\mathcal H]-1)^{D_r+1}=0$ for the tautological line bundle
$\mathcal H$ on $\mathbb{CP}^{D_r}$ therefore implies
$z^{D_r+1}=0$; see \cite[Corollary 2.7.9]{Atiyah67}.
Consequently \eqref{top:line-root} gives
\begin{equation}
 [p_{Y,j}^*\mathcal L_j]-1
 =(1+z)^{2^j}-1
 =\sum_{h=1}^{D_r}\binom{2^j}{h}z^h.
 \label{top:binomial}
\end{equation}
For $1\leq h\leq D_r$, let $v_2(h)$ be the exponent of $2$ dividing
$h$. The identity
\[
 h\binom{2^j}{h}=2^j\binom{2^j-1}{h-1}
\]
shows that $\binom{2^j}{h}$ is divisible by
$2^{j-v_2(h)}$, hence by
$2^{j-\lfloor\log_2D_r\rfloor}$. Every coefficient in
\eqref{top:binomial} has this common divisor. Multiplying the
resulting divisibility in $K^0$ by $p_{Y,j}^*e_3$ proves the claim.
\end{proof}

We have now obtained the growing divisor over the middle. It remains
to show that an equality between the classes at the two ends would
force their common pairing to have this divisor.

\subsection{Mayer--Vietoris detection and the direct limit}

For $r\geq3$, let
\[
 t_{\pm,j,r}:T_j\longrightarrow P_r(X_j)
\]
denote the inclusion of the triangulated torus at height $\pm3j$.

\begin{lemma}
\label{top:component-detection}
Fix an integer $r\geq3$ and a nonzero integer $n$. For all sufficiently
large $j$,
\begin{equation}
 n(t_{+,j,r})_*\tau_j-n(t_{-,j,r})_*\tau_j
 \neq0\quad\text{in }K_1(P_r(X_j)).
 \label{top:component-nonzero}
\end{equation}
\end{lemma}

\begin{proof}
In addition to the hypotheses needed for
Lemma~\ref{top:k-divisibility}, take $j>2r$. Cover $P_r(X_j)$ by the
induced subcomplexes
\[
 M^+=P_r(X_j)|_{\{t\geq-r\}},\qquad
 M^-=P_r(X_j)|_{\{t\leq r\}},
\]
and put $I=M^+\cap M^-$. These subcomplexes cover: a simplex not
contained in either would contain vertices at heights below $-r$ and
above $r$, contradicting its diameter bound because height is
$1$-Lipschitz. Their intersection is induced on the vertices at heights
in $[-r,r]$.

There is a simplicial map
\[
 h:I\longrightarrow P_r(Y_j)
\]
that forgets height. To check this, join any two vertices in a simplex
of $I$ by a graph path of length at most $r$ in $X_j$. Every vertex of
this path has height in $[-2r,2r]$, which lies in the product collar
since $j>2r$. Forgetting height sends the path to a path in $Y_j$ of
no greater length. Thus the images of the simplex vertices have
diameter at most $r$. The interpolation maps satisfy
\begin{equation}
 p_{X,j}|_I=p_{Y,j}\circ h.
 \label{top:collar-factorization}
\end{equation}

Let $a_\pm:I\hookrightarrow M^\pm$ and
$b_\pm:M^\pm\hookrightarrow P_r(X_j)$ be the inclusions, and write
$t^\pm:T_j\to M^\pm$ for the endpoint inclusions. The relevant part of
the Mayer--Vietoris sequence is
\[
 K_1(I)\xrightarrow{((a_+)_*,-(a_-)_*)}
 K_1(M^+)\oplus K_1(M^-)
 \xrightarrow{(b_+)_*+(b_-)_*}K_1(P_r(X_j)).
\]
If the difference in \eqref{top:component-nonzero} were zero,
exactness would give $\eta\in K_1(I)$ with
\[
 (a_+)_*\eta=n(t^+)_*\tau_j,
 \qquad (a_-)_*\eta=n(t^-)_*\tau_j.
\]
Pair the first equality with $(p_{X,j}|_{M^+})^*\theta_j$. Since
$p_{X,j}\circ t_{+,j,r}=\id_{T_j}$, naturality
\eqref{top:pairing-naturality} and \eqref{top:tau} give
\begin{equation}
 n
 =\langle(p_{X,j}|_I)^*\theta_j,\eta\rangle
 =\langle h^*p_{Y,j}^*\theta_j,\eta\rangle.
 \label{top:divisibility-contradiction}
\end{equation}
By Lemma~\ref{top:k-divisibility}, the last integer is divisible by
$2^{j-\lfloor\log_2D_r\rfloor}$. Taking $j$ large enough that
\[
 j-\lfloor\log_2D_r\rfloor>v_2(|n|)
\]
contradicts \eqref{top:divisibility-contradiction}.
\end{proof}

\begin{proposition}
\label{top:infinite-order}
The class $\alpha\in\KX_1(X)$ defined in \eqref{top:alpha} has infinite
order. In particular, $\alpha\otimes1\neq0$ in
$\KX_1(X)\otimes\mathbb Q$.
\end{proposition}

\begin{proof}
Fix a nonzero integer $n$ and an integer $r\geq3$. Let $\alpha_r$ be
the image of $\alpha_3$ under the map induced by
$P_3(X)\hookrightarrow P_r(X)$. For all sufficiently large $j$, the
component separation ensures that no simplex of $P_r(X)$ meets both
$X_j$ and its complement. Thus $P_r(X_j)$ is a clopen subcomplex, and
extension by zero on functions gives a $*$-homomorphism
\[
 C(P_r(X_j))\longrightarrow C_0(P_r(X)).
\]
Contravariance of $\KK$ in its first variable defines restriction
\[
 \KK_1(C_0(P_r(X)),\C)\longrightarrow K_1(P_r(X_j)).
\]
By \eqref{top:kk-product} and functoriality, the restriction of
$n\alpha_r$ is exactly the difference in
\eqref{top:component-nonzero}. Choosing $j$ also large enough for
Lemma~\ref{top:component-detection} proves $n\alpha_r\neq0$.

We have proved that for every nonzero $n$ and every integer $r\geq3$,
the scale-$r$ image of $n\alpha_3$ is nonzero. If $n\alpha$ vanished
in the filtered direct limit defining $\KX_1(X)$, it would vanish at
some finite later scale; integer scales are cofinal, so this is
impossible. Hence $\alpha$ has infinite order. The kernel of the map
from any abelian group to its tensor product with $\mathbb Q$ is its
torsion subgroup, which proves the final assertion.
\end{proof}
